\section{Kết luận}
\begin{enumerate}
    \item []\hspace*{0.7cm}Trong thời đại công nghệ số, việc thu thập và phân tích dữ liệu là yếu tố cốt lõi giúp các doanh nghiệp đưa ra quyết định kinh doanh hiệu quả. Thông qua đề tài này, nhóm đã nghiên cứu và ứng dụng công cụ mã nguồn mở Scrapy để tự động thu thập và phân tích dữ liệu từ trang web bán hàng trực tuyến. Kết quả thực nghiệm đã cho thấy hiệu quả của Scrapy trong việc thu thập thông tin một cách nhanh chóng, chính xác và tiết kiệm nguồn lực.
    
    \item[]Kết quả thực nghiệm cho thấy Selenium là một công cụ mạnh mẽ, cho phép thu thập dữ liệu từ các trang web động một cách nhanh chóng và chính xác, đồng thời tiết kiệm đáng kể nguồn lực. Dưới đây là những kết luận chính của đề tài:
    \begin{itemize}[left = 0.5cm]
        \item [$\bullet$]\textbf{Khả năng thu thập dữ liệu tự động:}Selenium đã chứng tỏ được tính hiệu quả trong việc tự động hóa quy trình thu thập dữ liệu từ các trang web có cấu trúc động, nơi mà nội dung có thể thay đổi theo thời gian. Việc sử dụng Selenium giúp giảm thiểu thời gian và công sức so với các phương pháp thu thập dữ liệu thủ công.
        
        \item [$\bullet$]\textbf{Phân tích xu hướng thị trường:} Dữ liệu thu thập từ trang web chứng khoán Simplize đã được xử lý và phân tích thành công. Nhóm nghiên cứu đã xác định được các xu hướng quan trọng trong thị trường chứng khoán dựa trên các chỉ số tài chính, biến động giá cổ phiếu và đánh giá từ người dùng. Kết quả này cung cấp cái nhìn sâu sắc về các cổ phiếu tiềm năng, từ đó hỗ trợ doanh nghiệp trong việc điều chỉnh chiến lược đầu tư.

        \item [$\bullet$]\textbf{Ứng dụng thực tiễn:}Dữ liệu thu thập được từ Selenium có thể được áp dụng trong nhiều lĩnh vực khác nhau, bao gồm phân tích tài chính, dự đoán xu hướng thị trường và phát triển các chiến lược đầu tư hiệu quả. Công cụ này không chỉ giúp các doanh nghiệp tiết kiệm chi phí mà còn nâng cao khả năng cạnh tranh trên thị trường tài chính.
    \end{itemize}
    \item[]Mặc dù Selenium mang lại nhiều lợi ích, nhưng nó cũng có một số hạn chế khi phải xử lý các trang web có cơ chế bảo mật phức tạp hoặc yêu cầu duy trì phiên đăng nhập lâu dài. Để giải quyết vấn đề này, cần kết hợp sử dụng với các công cụ khác như Beautiful Soup hoặc các phương pháp tối ưu hóa quy trình thu thập dữ liệu.

 
\end{enumerate}


\section{Kiến nghị}
\begin{enumerate}
    \item []\hspace*{0.7cm}Từ những kết quả từ quá trình thực nghiệm và phân tích, chúng tôi đưa ra một số kiến nghị để nâng cao hiệu quả ứng dụng của Selenium trong việc thu thập dữ liệu như:

    \begin{itemize}[left=0.3cm]
        \item[] \textbf{1. Kết hợp với Scrapy hoặc BeatufulSoup:} Đối với các trang web động sử dụng JavaScript, nên tích hợp Selenium với Scrapy hoặc BeautifulSoup để thu thập dữ liệu một cách hiệu quả hơn, mở rộng khả năng thu thập từ các trang web động.
        
        \item[]\textbf{2. Tối ưu hóa hiệu suất:} Cần tối ưu hóa mã và quy trình thu thập của Selenium, bao gồm việc cải thiện thời gian phản hồi và giảm thiểu lỗi để đảm bảo tốc độ và độ chính xác trong quá trình thu thập.

        \item[] \textbf{3. Phát triển mô hình dự báo:}Sử dụng dữ liệu thu thập được để xây dựng các mô hình dự báo xu hướng tiêu dùng, từ đó hỗ trợ doanh nghiệp trong việc ra quyết định nhanh chóng và hiệu quả.

        \item[] \textbf{4. Nâng cao khả năng quản lý dữ liệu lớn:} Khuyến khích sử dụng cơ sở dữ liệu phi cấu trúc như MongoDB để quản lý dữ liệu thu thập và áp dụng các công cụ phân tích mạnh mẽ như Pandas và NumPy để xử lý dữ liệu một cách nhanh chóng.

        \item[] \textbf{5. Đào tạo nhân lực:} Đầu tư vào đào tạo nhân viên về lập trình Python và hiểu biết về cách thức hoạt động của các công cụ thu thập dữ liệu mã nguồn mở nhằm nâng cao khả năng cạnh tranh và phát triển bền vững trong lĩnh vực này.
    \end{itemize}
\end{enumerate}